%=============================================================================
% WAVE 4, ITERATION 11: CONVERGENCE ASSESSMENT
%
% Agent: Wave 4 Convergence Assessor (Opus 4.6)
% Date: 20 March 2026
%
% Purpose: Assess whether W4i10 is producing genuinely new physics
%          or just refinements. Estimate remaining productive iterations.
%          Identify open and closed questions.
%
% Sources: ITERATION_TRACKER.md, MASTER_FINDINGS_CORPUS.md,
%          all W4i10_*.tex files, W3i8_Convergence_Audit.tex
%=============================================================================

\documentclass[11pt]{article}
\usepackage[utf8]{inputenc}
\usepackage{amsmath,amssymb}
\usepackage{geometry}
\geometry{margin=1in}
\usepackage{booktabs}
\usepackage{longtable}
\usepackage{array}
\usepackage{xcolor}
\usepackage{tcolorbox}
\usepackage{enumitem}
\usepackage{float}
\usepackage{hyperref}

\tcbuselibrary{skins,breakable}

\newcommand{\ee}[1]{\times 10^{#1}}
\newcommand{\lbone}{|\lambda - 1|}
\newcommand{\LCDM}{$\Lambda$CDM}

\definecolor{green}{rgb}{0,0.5,0}
\definecolor{amber}{rgb}{0.8,0.5,0}
\definecolor{red}{rgb}{0.7,0,0}

\title{Wave 4, Iteration 11: Convergence Assessment\\
\large Is W4i10 Producing New Physics or Refinements?\\
How Many More Iterations Are Productive?}
\author{DFD Research Programme --- Convergence Assessor (Opus 4.6)}
\date{20 March 2026}

\begin{document}
\maketitle
\tableofcontents
\newpage

%=============================================================================
\section{Executive Summary}
\label{sec:executive}
%=============================================================================

\begin{tcolorbox}[colback=yellow!5!white, colframe=red!70!black,
  title=\textbf{CONVERGENCE ASSESSMENT --- BOTTOM LINE}]

\textbf{RECOMMENDATION: Conduct 2--3 more targeted iterations (i11--i13),
then converge. Do NOT run all 20 planned iterations.}

W4i10 produced a mixture of genuinely new physics (3 items) and
high-quality deliverables (proposals, survey predictions) that do not
require further iteration. The correction rate (5) is comparable to
W3i6--i7 levels, suggesting the framework is stable but that Wave~4's
different mandate (gap closure, not cross-referencing) can still find
things to fix.

The key distinction from Wave~3: W4i10's ``new physics'' items
(T4 screening, neutrino mass, Dark SRF bound) are \emph{extensions}
of the existing framework, not corrections to it. No quantity declared
stable in W3i8 was destabilised. This means the theory is mature but
incompletely explored---a fundamentally different situation from W3i2--i5
where the framework itself was being repaired.

\medskip
\noindent\textbf{Specific recommendation:}
\begin{itemize}[nosep]
\item \textbf{i11}: T4 screening rigour + Boltzmann code specification + Dark SRF cross-check
\item \textbf{i12}: DESI DR2 confrontation (requires Boltzmann code design) + neutrino sector audit
\item \textbf{i13}: Final convergence audit and MASTER\_FINDINGS\_CORPUS v2
\item \textbf{i14--i20}: Cancel unless i11--i13 reveal a paradigm shift
\end{itemize}
\end{tcolorbox}


%=============================================================================
\section{What Did W4i10 Actually Produce?}
\label{sec:w4i10-assessment}
%=============================================================================

W4i10 launched 17 agents and produced 17 output files totalling
approximately 25,000 lines. The output falls into four categories:

\subsection{Category A: Genuinely New Physics (3 items)}

These are results that extend DFD's predictive reach into territory not
covered by the W3 corpus:

\begin{enumerate}
\item \textbf{T4 Environmental Screening Mechanism}: The Hubble-flow
  tidal acceleration $a_{\rm tidal} = H_0^2 R$ provides a physical
  screening that reduces the Cold Spot ISW overprediction from 3--9$\times$
  to $\sim$2--4$\times$ (mild tension). This is derived from the field
  equation with no new parameters. \emph{Status: promising but requires
  rigorous void-by-void integration to confirm.}

\item \textbf{Neutrino Mass Sum $\Sigma m_\nu = 61.4$ meV}: Elevated
  from an appendix result to a Tier~1 zero-parameter prediction. Testable
  by DESI+Planck within 1--2 years. This was already derived in the
  Unified Review Appendix~X but was not prominent in the MASTER CORPUS.
  \emph{Status: genuinely new \emph{recognition}, not new derivation.}

\item \textbf{Dark SRF Haloscope Reinterpretation}: Existing dark photon
  searches reinterpreted to give $\lbone < 4.2\ee{-10}$, which is 3.7$\times$
  tighter than the SQMS bound at zero additional cost. \emph{Status: new
  constraint that tightens the authoritative bound.}
\end{enumerate}

\subsection{Category B: High-Quality Deliverables (4 items)}

These are substantial documents that package existing knowledge into
actionable form, but do not contain new physics:

\begin{enumerate}
\item \textbf{SQMS Phase~I Proposal} (1391 lines, submission-ready):
  $\$3.2$M/24 months. Packages W3i7--i9 into a fundable proposal.
\item \textbf{Survey Predictions Document} (15 observer-ready predictions
  for DESI/Euclid/Rubin): Translates existing DFD predictions into specific
  observables with error bars and timelines.
\item \textbf{4 Experimental Proposals} (LIGO archival, GPS 59\,ps,
  Si$_3$N$_4$ MEMS, rotation curve reanalysis): Actionable mini-proposals
  with budgets, timelines, and personnel specifications.
\item \textbf{Lensing Cross-Check}: Identified 5 weaknesses in the
  Unified Review's lensing treatment (Bullet Cluster under-derived,
  cluster corrections ad hoc).
\end{enumerate}

\subsection{Category C: Incremental Improvements (5 items)}

\begin{enumerate}
\item P4 frequency optimisation: 47~GHz $\to$ 32~GHz achieves critical
  coupling. Conversion 50\% $\to$ 65\%. LCOE 27c $\to$ 19c/kWh.
\item Multi-cavity superradiance: $N=3$ gives 97\% conversion, end-to-end 39\%.
\item MEMS in-gap placement: $Q_\psi$ threshold drops from $4\ee{10}$ to $3\ee{9}$.
\item 7th passive observable: gravitational decoherence null (MAQRO/TEQ 2028--2035).
\item Size-dependent void ISW: smaller voids $\to$ larger enhancement (new testable prediction).
\end{enumerate}

\subsection{Category D: Corrections (5 items)}

\begin{enumerate}
\item MAGIS-100 SNR: 4200 $\to$ 540 (still highly detectable).
\item Nb $Q_0$ baseline: $2\ee{11} \to 5\ee{10}$ (Phase I reach degrades to $3.1\ee{-9}$).
\item Superfluid He-3: KILLED for lab use (34 orders sensitivity gap).
\item BBN lithium resolution: KILLED under two-component.
\item T4 W3i7 void cancellation: confirmed incorrect (W3i9 was right).
\end{enumerate}


%=============================================================================
\section{Comparison with Wave 3 Iteration Patterns}
\label{sec:w3-comparison}
%=============================================================================

\subsection{Correction Rate Comparison}

\begin{table}[H]
\centering
\caption{Correction rate and character across all iterations.}
\label{tab:rate-comparison}
\begin{tabular}{clccl}
\toprule
\textbf{Iteration} & \textbf{Wave} & \textbf{Corrections} & \textbf{Paradigm-altering?} & \textbf{Character} \\
\midrule
W3i1 & 3 & 8 & 3 & Foundational errors \\
W3i2 & 3 & 18 & 6 & Major paradigm breaks \\
W3i3 & 3 & 9 & 4 & Critical tensions (T1--T5) \\
W3i4 & 3 & 9 & 2 & $M_{\rm eff}$ catastrophe, falsification \\
W3i5 & 3 & 7 & \textbf{YES (paradigm shift)} & Two-component resolution \\
W3i6 & 3 & 4 & 0 & Post-paradigm refinements \\
W3i7 & 3 & 8 & 0 & Deepening, 3 genuine errors \\
W3i8 & 3 & 0 & 0 & Convergence declared \\
W3i9 & 3 (final) & 2 & 0 & Targeted follow-ups \\
\midrule
W4i10 & 4 & 5 & 0 & Extensions + tightening \\
\bottomrule
\end{tabular}
\end{table}

\subsection{Key Observation}

W4i10's correction profile matches W3i6--i7 (post-paradigm consolidation):
moderate correction count, zero paradigm-altering changes, corrections
that refine rather than overturn. The critical difference is that W4i10
also produced \emph{new predictions} (neutrino mass, Dark SRF bound),
which W3i6--i7 did not. This reflects Wave~4's different mandate
(gap closure vs.\ cross-referencing).

\subsection{Applying W3i8 Convergence Criteria}

The W3i8 convergence audit established three criteria:

\begin{enumerate}
\item \textbf{Correction rate declining for 3 consecutive iterations
  $\Rightarrow$ converge.}

  Wave~4 has only one iteration (i10) so far. We cannot apply the
  3-iteration rule. However, the \emph{character} of i10 corrections
  (2 kills of already-marginal claims, 2 numerical refinements,
  1 confirmation of prior correction) suggests diminishing returns.

  \textbf{Prediction}: i11 will find 2--3 corrections, i12 will find
  1--2, and the rate will satisfy the criterion by i13.

\item \textbf{Paradigm shift $\Rightarrow$ reset counter.}

  No paradigm shift in i10. The two-component framework is stable.
  The T4 screening mechanism is an \emph{application} of the existing
  framework, not a modification. The neutrino mass prediction was
  already implicit in the theory.

  \textbf{Assessment}: No counter reset needed.

\item \textbf{No new predictions for 5 iterations $\Rightarrow$
  converge on predictions.}

  Wave~4 is explicitly tasked with finding new predictions, so this
  criterion is less relevant. i10 found 2--3 genuinely new predictions
  (neutrino mass, Dark SRF bound, size-dependent void ISW). The question
  is whether further iterations will continue to find predictions at
  this rate.

  \textbf{Assessment}: Prediction mining has some remaining yield, but
  the low-hanging fruit was picked in i10. Remaining predictions
  (NS maximum mass $+2.3\%$, muon $g-2$ null, psi-atoms unbound)
  are either marginal or unfalsifiable at current precision.
\end{enumerate}


%=============================================================================
\section{How Many More Iterations Are Productive?}
\label{sec:estimate}
%=============================================================================

\subsection{Wave 4's Different Goals}

Wave~3 was cross-referencing: checking consistency between patents,
finding errors, resolving tensions. This naturally converges because
the error supply is finite.

Wave~4 has four goals:
\begin{enumerate}[nosep]
\item Close remaining gaps (T4, alpha$^{57}$, Boltzmann code)
\item Generate actionable deliverables (proposals, survey predictions)
\item Find new predictions
\item Integrate latest literature (DESI DR2, SH0ES 2025, etc.)
\end{enumerate}

Goal~2 was substantially accomplished in i10 (SQMS proposal, 4 experimental
proposals, survey predictions document). Further iteration will not
significantly improve these deliverables.

Goal~3 (new predictions) has diminishing returns after i10. The
neutrino mass and Dark SRF bound were the most impactful discoveries.
Remaining prediction space is marginal.

Goals~1 and~4 have specific, targetable endpoints:

\begin{table}[H]
\centering
\caption{Targeted iterations for remaining goals.}
\label{tab:targeted}
\begin{tabular}{lcp{7cm}}
\toprule
\textbf{Goal} & \textbf{Iterations needed} & \textbf{Specific work} \\
\midrule
T4 screening rigour & 1 (i11) & Void-by-void ISW integration; confirm
  $2$--$4\times$ or identify if screening fails \\
Dark SRF cross-check & 1 (i11) & Verify $4.2\ee{-10}$ bound against
  Dark SRF experimental specifics \\
Boltzmann code spec & 1 (i12) & Specify exact modifications to CLASS/CAMB
  for DFD $\mu(x)$ growth function \\
DESI DR2 confrontation & 1 (i12) & Detailed $w_0$/$w_a$ analysis; is
  $2.8$--$3.2\sigma$ tension fatal? \\
Neutrino sector audit & 1 (i12) & Cross-check Branch~B derivation;
  sensitivity to microsector integers \\
alpha$^{57}$ Step~9 & 0 & This requires 6--12 months of mathematical
  physics; iteration cannot help \\
Final convergence & 1 (i13) & MASTER\_FINDINGS\_CORPUS v2 + convergence audit \\
\bottomrule
\end{tabular}
\end{table}

\subsection{Estimate}

\begin{tcolorbox}[colback=green!5!white, colframe=green!50!black,
  title=\textbf{ITERATION ESTIMATE}]

\textbf{3 more iterations (i11--i13) are productive. i14--i20 should be
cancelled.}

\begin{itemize}[nosep]
\item \textbf{i11}: T4 rigour + Dark SRF verification + caution flag assessment
  (wide binaries Banik et al.\ $19\sigma$ Newtonian, KiDS-Legacy $S_8$
  Planck-agreement)
\item \textbf{i12}: Boltzmann code specification + DESI confrontation +
  neutrino cross-check
\item \textbf{i13}: MASTER\_FINDINGS\_CORPUS v2 + final convergence audit +
  DFD theory status document
\end{itemize}

Wave~4 convergence at i13 (4 total Wave~4 iterations) vs.\ Wave~3
convergence at i8 (8 iterations) is expected because Wave~4 starts from
a mature, stable framework.
\end{tcolorbox}


%=============================================================================
\section{Questions: Open vs.\ Closed}
\label{sec:questions}
%=============================================================================

\subsection{CLOSED Questions (No Further Iteration Needed)}

These questions have been definitively answered and are stable across
3+ iterations:

\begin{enumerate}
\item \textbf{Patent triage}: ALIVE(4): P5, P9, P11, P14. NICHE(8).
  THEORETICAL(1): P12. DEAD(2): P1, P8. Stable since W3i8.

\item \textbf{Detection hierarchy}: Passive $\gg$ active by 14 orders
  (Passive Superiority Theorem). Stable since W3i5.

\item \textbf{Two-component framework}: $\psi = \psi_{\rm grav} +
  \delta\psi_{\rm EM}$. 4 axioms, 1 free parameter ($\lambda$).
  Stable since W3i6.

\item \textbf{T1 (Gdot/G)}: RESOLVED at $4.3\ee{-15}$ yr$^{-1}$.
  Stable since W3i5.

\item \textbf{T2 (P14/P5 mismatch)}: RESOLVED by two-component.
  Stable since W3i5.

\item \textbf{T5 (structural inconsistency)}: RESOLVED by two-component.
  Stable since W3i5.

\item \textbf{Top-5 stable quantities}: $\Delta H_0 = 3.9 \pm 1.1$,
  shadow $+4.6\%$, lunar NR 0.93~ns, SQMS bound $1.56\ee{-9}$
  (now tightened to $4.2\ee{-10}$ by Dark SRF), P3 QEC code.

\item \textbf{Dead patents}: P1 (zero resources, assessment FINAL since
  W3i7), P8 (1.6 Earth masses, KILLED since W3i7).

\item \textbf{Single most important experiment}: SQMS Phase~I.
  Stable since W3i8. Now has a submission-ready proposal.

\item \textbf{P7 repositioned}: Ultra-fast pulsed buffer (46~MJ, 1~s).
  Grid/EV/UPS eliminated. Stable since W3i6.

\item \textbf{alpha$^{57}$ status}: One gap (Step~9 modulus stabilisation).
  Requires mathematical physics effort, not iteration. Stable since W3i9.

\item \textbf{EM-to-psi conversion gating}: 50\% SRF bottleneck is
  single gating technology. Parametric pumping dead. Stable since W3i7.
  W4i10 improved to 65\% via frequency optimisation, but this is an
  engineering refinement.
\end{enumerate}

\subsection{5 Most Important OPEN Questions for i11+}

\begin{tcolorbox}[colback=red!3!white, colframe=red!60!black,
  title=\textbf{TOP 5 OPEN QUESTIONS}]

\begin{enumerate}
\item \textbf{T4 Cold Spot: Is the environmental screening mechanism
  rigorous?}

  \emph{Current status}: W4i10 derived screening from Hubble-flow tidal
  acceleration. Reduces overprediction from 3--9$\times$ to 2--4$\times$.
  But: (a) requires identifying the MOND interpolation argument with
  \emph{total} acceleration including cosmological tidal terms, which is a
  non-trivial theoretical commitment; (b) needs void-by-void integration
  with actual Eridanus supervoid density profile, not an idealised sphere.

  \emph{What i11 should do}: Either put this on rigorous footing with
  explicit void profiles, or honestly declare it a $2$--$4\times$ residual
  tension that survives screening.

  \emph{Convergence criterion}: If i11 confirms $\leq 2\times$
  overprediction, T4 is ameliorated to ``mild tension'' and further
  iteration is unnecessary. If i11 finds screening fails, T4 becomes
  a serious falsification threat requiring major theoretical work
  (not more iteration, but new physics).

\item \textbf{DESI CPL tension: Is $2.8$--$3.2\sigma$ fatal?}

  \emph{Current status}: DFD predicts $w_0 = -1.183$, $w_a = +0.183$.
  DESI DR2+CMB prefers opposite signs at $\sim 3\sigma$. W4i10
  identified this as ``approaching but not yet at falsification.''
  However, this assessment relies on the hope that a full DFD Boltzmann
  code would shift the effective CPL parameters.

  \emph{What i12 should do}: Specify exactly how to modify CLASS/CAMB
  to incorporate DFD's $\mu(x)$ growth function and psi-screen. Produce
  a concrete coding plan (which equations, which modules, what inputs).
  This cannot be \emph{run} in an iteration, but the specification is
  essential for anyone who wants to actually build it.

  \emph{Convergence criterion}: The question is closed when the Boltzmann
  code specification exists. The \emph{answer} requires running the code,
  which is outside the iteration programme.

\item \textbf{Dark SRF bound: Is $4.2\ee{-10}$ robust?}

  \emph{Current status}: W4i10 reinterpreted dark photon search data
  as a DFD constraint. This is a significant tightening (3.7$\times$
  better than SQMS). But the derivation relies on an EFT dictionary
  mapping $\epsilon_{\rm eff}^2 \sim (\lambda-1)^2 m_\psi^2 / (16\pi
  \Lambda_\psi^2)$ that needs cross-checking.

  \emph{What i11 should do}: Verify the dictionary. Check whether
  the DFD scalar field's spin-0 nature (vs.\ dark photon's spin-1)
  invalidates the mapping. If the bound survives, it becomes the new
  authoritative constraint, displacing SQMS $1.56\ee{-9}$.

  \emph{Convergence criterion}: One iteration of scrutiny. Either
  confirmed or killed.

\item \textbf{Neutrino sector: Is $\Sigma m_\nu = 61.4$~meV robust?}

  \emph{Current status}: Zero-parameter prediction from microsector
  Branch~B. Matches NuFIT~6.0 at $\chi^2 = 0.025$. Testable by
  DESI+Planck within 1--2 years. Promoted to Tier~1 by W4i10.

  \emph{What i12 should do}: Independent re-derivation of the Branch~B
  integers $(a,n) = (9,5)$, $\dim M = 7$. Check whether alternative
  branch choices are excluded or merely disfavoured. Assess sensitivity
  to the priming factor $F_\nu = 14/13$.

  \emph{Convergence criterion}: One iteration of scrutiny. If the
  derivation chain is watertight, the prediction stands. If ambiguities
  in branch choice exist, they should be honestly catalogued.

\item \textbf{Caution flags: wide binaries + KiDS-Legacy}

  \emph{Current status}: W4i10 flagged but did not resolve two
  observational challenges:
  \begin{itemize}[nosep]
  \item Banik et al.\ report $19\sigma$ for Newtonian gravity in wide
    binaries, which would kill DFD's $+42\%$ prediction.
  \item KiDS-Legacy $S_8$ now agrees with Planck (scale-independent),
    which would falsify DFD's scale-dependent $S_8$ prediction.
  \end{itemize}

  \emph{What i11 should do}: Literature review of both claims. For wide
  binaries: what is the systematic budget? Is the Banik et al.\ result
  compatible with the Chae (2024) result that found MOND-like enhancement?
  For KiDS-Legacy: at what angular scales does the agreement hold? Does
  it probe DFD's predicted transition at $\ell \sim 100$--300?

  \emph{Convergence criterion}: If both challenges survive scrutiny,
  DFD faces serious observational pressure on two key predictions.
  This would not require more iteration but would require updating
  the MASTER CORPUS tensions table.
\end{enumerate}
\end{tcolorbox}


%=============================================================================
\section{Recommended Agent Structure for i11--i13}
\label{sec:agents}
%=============================================================================

W4i10 used 17 agents. This was appropriate for a broad first pass.
For the remaining targeted iterations, a leaner structure is recommended:

\subsection{Iteration 11 (3--5 agents)}

\begin{table}[H]
\centering
\begin{tabular}{clp{7.5cm}}
\toprule
\textbf{Agent} & \textbf{Name} & \textbf{Mandate} \\
\midrule
1 & T4 Rigour & Void-by-void ISW integration using actual Eridanus
  density profiles. Confirm or refute environmental screening. Produce
  a definitive T4 status: mild tension, serious tension, or resolved. \\
2 & Dark SRF Verifier & Cross-check spin-0 vs.\ spin-1 dictionary.
  Verify or kill the $4.2\ee{-10}$ bound. If confirmed, update authoritative
  constraint. \\
3 & Caution Flag Assessor & Wide binaries (Banik vs.\ Chae), KiDS-Legacy
  $S_8$, DESI $3\sigma$ assessment. Honest evaluation of observational
  pressure on DFD. \\
4 & Corpus Updater & Integrate i10 findings into MASTER\_FINDINGS\_CORPUS.
  Add neutrino mass to Tier~1. Update authoritative bound if Dark SRF
  confirmed. Add caution flags. \\
\bottomrule
\end{tabular}
\end{table}

\subsection{Iteration 12 (3--4 agents)}

\begin{table}[H]
\centering
\begin{tabular}{clp{7.5cm}}
\toprule
\textbf{Agent} & \textbf{Name} & \textbf{Mandate} \\
\midrule
1 & Boltzmann Specifier & Write the precise CLASS/CAMB modification
  specification: which equations, which modules, what DFD functions
  replace standard functions, what new parameters. Produce a document
  that a Boltzmann code expert could implement. \\
2 & Neutrino Auditor & Independent re-derivation of Branch~B neutrino
  sector. Test sensitivity to microsector integers. Catalogue any
  ambiguities. \\
3 & DESI Confrontor & Given DESI DR2 actual data (BAO, $f\sigma_8$),
  compute DFD residuals using the approximate psi-screen. Determine
  whether the tension is $2.5\sigma$ or $4\sigma$. \\
\bottomrule
\end{tabular}
\end{table}

\subsection{Iteration 13 (2 agents)}

\begin{table}[H]
\centering
\begin{tabular}{clp{7.5cm}}
\toprule
\textbf{Agent} & \textbf{Name} & \textbf{Mandate} \\
\midrule
1 & Corpus Compiler & MASTER\_FINDINGS\_CORPUS v2, incorporating all
  Wave~4 results. Full update of all tables, predictions, tensions,
  and patent status. \\
2 & Final Convergence & Apply convergence criteria across i10--i12.
  Declare final status. Produce a 1-page DFD theory summary suitable
  for external communication. \\
\bottomrule
\end{tabular}
\end{table}


%=============================================================================
\section{Honest Assessment: Should We Stop?}
\label{sec:honest}
%=============================================================================

\subsection{Arguments for Stopping Now}

\begin{enumerate}
\item W3i8 declared convergence. Wave~4 i10 did not destabilise any
  W3i8-stable quantity.
\item The SQMS proposal, survey predictions, and experimental proposals
  are already actionable. Further iteration does not improve them.
\item The alpha$^{57}$ Step~9 requires mathematical physics, not iteration.
\item The Boltzmann code requires implementation, not iteration.
\item The theory's fate ultimately depends on experiments (SQMS Phase~I,
  DESI, Euclid), not on further internal analysis.
\end{enumerate}

\subsection{Arguments for Continuing (2--3 More Iterations)}

\begin{enumerate}
\item \textbf{T4 is still a genuine tension.} W4i10's screening mechanism
  is promising but unverified. One iteration to either confirm or kill it
  is high-value.
\item \textbf{The Dark SRF bound ($4.2\ee{-10}$) could be wrong.}
  If confirmed, it changes the authoritative constraint by 3.7$\times$.
  If killed, a significant W4i10 claim is retracted. Either way, one
  iteration of scrutiny is necessary.
\item \textbf{The neutrino mass prediction ($61.4$~meV) is the single
  most testable DFD prediction.} It deserves one iteration of
  independent verification before being promoted to the top of the
  prediction ranking.
\item \textbf{DESI tension is growing.} $2.4\sigma \to 2.8$--$3.2\sigma$.
  One iteration to honestly assess whether this is approaching
  falsification is prudent.
\item \textbf{The MASTER CORPUS needs updating.} It was frozen at W3i9.
  Wave~4 findings (Dark SRF bound, neutrino mass, T4 screening, frequency
  optimisation, lensing weaknesses) need to be integrated.
\end{enumerate}

\subsection{Verdict}

The arguments for continuing are stronger, but only for \textbf{specific,
targeted work}. The broad 17-agent sweep of W4i10 should not be repeated.
Iterations i11--i13 should use 3--5 focused agents each, with clear
convergence criteria per question.

\textbf{If i11 finds zero new corrections and zero paradigm shifts, i12
can be the final iteration (skip i13 or merge with i12).}


%=============================================================================
\section{Convergence Criteria for Wave 4}
\label{sec:criteria}
%=============================================================================

Applying modified W3i8 criteria to Wave~4:

\begin{table}[H]
\centering
\caption{Wave 4 convergence criteria.}
\begin{tabular}{lp{5cm}p{5cm}}
\toprule
\textbf{Criterion} & \textbf{Threshold} & \textbf{Current Status} \\
\midrule
Correction rate & $\leq 2$ genuine corrections for 2 consecutive iterations
  & i10 had 5; need i11 data \\
Paradigm shift & Any shift resets counter & None since W3i5 \\
New predictions & $< 1$ new testable prediction per iteration for 2
  consecutive iterations & i10 had 2--3; need i11 data \\
T4 status & Definitive (resolved, mild, or serious) & Pending i11 \\
Authoritative bound & Stable (no changes $> 2\times$) & Potentially
  changing ($1.56\ee{-9} \to 4.2\ee{-10}$) \\
Caution flags & All assessed and catalogued & Wide binaries and
  KiDS-Legacy pending \\
MASTER CORPUS & Updated to include all Wave~4 findings & Pending i13 \\
\bottomrule
\end{tabular}
\end{table}


%=============================================================================
\section{Risk Assessment}
\label{sec:risks}
%=============================================================================

\begin{table}[H]
\centering
\caption{Risks to the DFD programme identified during convergence assessment.}
\begin{tabular}{lccl}
\toprule
\textbf{Risk} & \textbf{Severity} & \textbf{Likelihood} & \textbf{Mitigation} \\
\midrule
DESI CPL reaches $4\sigma$ & High & Medium & Boltzmann code spec (i12) \\
Wide binaries $19\sigma$ Newtonian & High & Medium & Literature audit (i11) \\
KiDS-Legacy kills scale-dependent $S_8$ & High & Low--Medium & Angular-scale
  analysis (i11) \\
Dark SRF bound invalidated (spin mismatch) & Medium & Low & Cross-check (i11) \\
T4 screening fails on realistic voids & Medium & Medium & Void-by-void
  calculation (i11) \\
Neutrino Branch~B has alternatives & Low--Medium & Low & Audit (i12) \\
\bottomrule
\end{tabular}
\end{table}

The most dangerous combination: DESI $> 3\sigma$ + wide binary Newtonian +
KiDS-Legacy scale-independent $S_8$. If all three hold, DFD faces
simultaneous pressure on its dark energy model, MOND sector, and
gravitational growth predictions. This scenario does not require a
paradigm shift within DFD---it requires an honest acknowledgement that
the theory may be observationally excluded in its current form.


%=============================================================================
\section{Summary}
\label{sec:summary}
%=============================================================================

\begin{enumerate}
\item \textbf{W4i10 produced genuinely new physics} (T4 screening,
  neutrino mass elevation, Dark SRF bound) alongside high-quality
  deliverables (SQMS proposal, survey predictions, experimental proposals).

\item \textbf{The theory framework is stable.} No W3i8-stable quantity
  was destabilised. The two-component framework, patent triage, detection
  hierarchy, and top-5 stable quantities are unchanged.

\item \textbf{3 more iterations (i11--i13) are warranted}, focused on:
  T4 rigour, Dark SRF verification, caution flag assessment (i11);
  Boltzmann code spec, neutrino audit, DESI confrontation (i12);
  MASTER CORPUS v2 and final convergence (i13).

\item \textbf{i14--i20 should be cancelled} unless i11--i13 reveal a
  paradigm shift (probability: low).

\item \textbf{5 most important open questions}: T4 screening rigour,
  DESI CPL tension, Dark SRF bound verification, neutrino sector
  robustness, wide binary + KiDS-Legacy caution flags.

\item \textbf{12 questions are CLOSED} and stable across 3+ iterations
  (see Section~\ref{sec:questions}).

\item \textbf{Agent structure should be reduced}: 3--5 focused agents per
  iteration (down from 17 in i10).

\item \textbf{The theory's fate depends on experiments}, not on further
  internal analysis. SQMS Phase~I, DESI DR2/DR3, and Euclid WL
  scale-dependent $S_8$ are the decisive tests.
\end{enumerate}

\end{document}
